\begin{textsample}{POS Dim 3 – 2021\_c5 – Score 2.00 – 2021\_c5/enade\_2021\_ef\_be61acbd52.txt}  \label{ex:f3_pos_083}
CAPÍTULO 5 DISTRIBUIÇÃO DOS CONCEITOS 5.1 PANORAMA NACIONAL DA DISTRIBUIÇÃO DOS CONCEITOS O Conceito Enade é um indicador de qualidade que avalia os cursos por intermédio dos desempenhos dos estudantes a partir dos resultados do Enade.
Os valores possíveis vão de 1 ( pior situação ) a 5 ( melhor situação ).
Os cursos com apenas um ou sem nenhum concluinte participante não obtêm o Conceito Enade, ficando Sem Conceito ( SC ).
Foram avaliados em Educação Física 1.147 cursos, sendo 542 ( 47,3 \% ) de Licenciatura e 605 ( 52,7 \% ) de Bacharelado.
Enfatiza -se, mais uma vez, que a diferença entre os cursos tabulados neste capítulo e no Capítulo 2 corresponde aos cursos que não foram avaliados, ou seja, àqueles, em princípio, sem alunos concluintes que tivessem sido inscritos para a prova.
Na Tabela 5.1a, são apresentadas a quantidade e a distribuição de cursos de Licenciatura de Educação Física participantes do Enade/2021, por faixa de conceito e Grande Região.
Observando -se os dados da Tabela 5.1a, nota -se que, dos 542 cursos participantes, 219 ( 40,4 \% ) classificaram -se com conceito 3, o valor modal.
Esse foi também o conceito modal de todas as regiões : Norte ( 46,7 \% ), Nordeste ( 41,6 \% ), Sudeste ( 39,0 \% ), Sul ( 40,4 \% ) e Centro-Oeste ( 41,2 \% ).
O conceito 2 foi o segundo mais frequente em nível nacional ( 27,7 \%, correspondendo a 150 cursos ).
Houve, ainda, 113 cursos ( 20,8 \% ) que receberam o conceito 4, 24 cursos ( 4,4 \% ) que receberam o conceito 5, e 13 cursos ( 2,4 \% ) que receberam o conceito 1.
Dos 542 cursos de Educação Física, 23 ( 4,2 \% ) ficaram Sem Conceito ( SC ).
A região Norte participou com 30 cursos, o correspondente a 5,5 \% do total nacional.
Desses, 14 cursos receberam o conceito 3, o que equivale a 46,7 \% do total regional, o conceito modal na região, como já comentado.
A nove cursos ( 30,0 \% ) atribuiu -se o conceito 2, a cinco cursos ( 16,7 \% ), o conceito 4, e a um ( 3,3 \% ), o conceito 5.
Nenhum curso recebeu conceito 1, e um curso ( 3,3 \% ) ficou Sem Conceito ( SC ).
A região Nordeste participou com 101 cursos, o equivalente a 18,6 \% do total nacional.
Nessa região, 42 cursos ( 41,6 \% em termos regionais ) obtiveram conceito 3, o conceito modal para a região, como já comentado.
O conceito 4 foi atribuído a 26 cursos ( 25,7 \% ) ; o conceito 5, a oito cursos ( 7,9 \% ) ; o conceito 1, a dois cursos ( 2,0 \% ) ; e o conceito 2, a 19 cursos ( 18,8 \% ).
Quatro cursos ( 4,0 \% ) ficaram Sem Conceito ( SC ).
Dos 251 cursos participantes da região Sudeste ( 46,3 \% do total nacional ), o conceito 3 foi o conceito modal, como já anteriormente assinalado, tendo sido obtido por 98 cursos ( 39,0 \% ).
O conceito 1 foi atribuído a quatro cursos ( 1,6 \% ) ; o conceito 2, a 78 cursos ( 31,1 \% ) ; o conceito 4 a 47 cursos ( 18,7 \% ) ; e o conceito 5, a 12 cursos ( 4,8 \% ).
Nessa região, 12 cursos ( 4,8 \% ) ficaram Sem Conceito ( SC ).
Os 109 cursos da região Sul corresponderam a 20,1 \% do total nacional.
Como já indicado, houve predominância do conceito 3, atribuído a 40,4 \% dos cursos da região ( conceito modal ), ou seja, atribuído a 44 dos 109 cursos participantes na região Sul.
O conceito 4 foi atribuído a 30 cursos ( 27,5 \% ), o conceito 5, a dois cursos ( 1,8 \% ) ; o conceito 1, a dois cursos também ( 1,8 \% ) ; e o conceito 2, a 27 cursos ( 24,8 \% ).
Quatro cursos ( 3,7 \% ) ficaram Sem Conceito ( SC ) na região Sul.
Já dos 51 cursos participantes na região Centro-Oeste ( 9,4 \% do total nacional ), 21 ( 41,2 \% em termos regionais ) receberam conceito 3, o conceito modal, como já destacado.
Os conceitos 1 e 4 foram atribuídos a cinco cursos cada ( 9,8 \% em cada conceito ), 17 cursos ( 33,3 \% ) obtiveram conceito 2, e um obteve ( 2,0 \% ) o conceito 5.
Nessa região, dois cursos ( 3,9 \% ) ficaram Sem Conceito ( SC ).
No Gráfico 5.1a, é apresentada a distribuição cumulativa do Conceito Enade dos cursos de Licenciatura segundo Grande Região.
Em um gráfico de distribuição cumulativa, em cada valor de abscissa ( eixo x ), o valor apresentado no eixo das ordenadas ( eixo y ) é o valor acumulado da variável ( conceito Enade ).
Observando -se a linha vertical pontilhada, que \textbf{ajuda} a visualizar o percentual de conceitos acumulados até o valor 3, ou seja, os conceitos 1, 2 e 3, pode -se notar que, por exemplo, a região Norte ( linha azul ) apresentou 79,3 \% dos cursos com conceito 3 ou abaixo desse valor.
A região Norte ( linha verde ), por outro lado, apresentou 64,9 \% dos cursos na mesma situação : com conceito 3 ou abaixo.
Em linhas gerais, em um gráfico de distribuição cumulativa, poligonais mais à esquerda representam as Grandes Regiões com a pior distribuição, e poligonais mais à direita, as Grandes Regiões com os melhores conceitos.
A região Centro-Oeste ( linha \textbf{laranja} ) apresentou os cursos com os conceitos mais baixos, ao passo que a região Nordeste ( linha verde ) apresentou a distribuição com os valores maiores.
Cursos em IES situadas nas regiões Norte ( linha azul ), Sudeste ( linha vermelha ) e Sul ( linha rosa ) apresentaram uma situação intermediária. verde ) apresentou a distribuição com os valores maiores.
Cursos em IES situadas nas regiões Norte ( linha azul ), Sudeste ( linha vermelha ) e Sul ( linha rosa ) apresentaram uma situação intermediária.
Na Tabela 5.1b, são apresentadas a quantidade e a distribuição de cursos de Bacharelado de Educação Física participantes do Enade/2021, por faixa de conceito e Grande Região.
Observando -se os dados da Tabela 5.1b, nota -se que, dos 605 cursos participantes, 272 ( 45,0 \% ) classificaram -se com conceito 3, o valor modal.
Esse foi também o conceito modal na maioria das regiões, exceto na Centro-Oeste.
O conceito 2 foi o segundo mais frequente em nível nacional ( 28,1 \%, correspondendo a 170 cursos ), tendo sido esse o conceito modal na região Centro-Oeste ; já o conceito 4 foi o terceiro mais frequente ( 17,2 \%, correspondendo a 104 cursos ).
O quarto conceito mais frequente foi o 5 ( 6,6 \% correspondendo a 40 cursos ) e, por último, o conceito 1 ( 2,1 \% correspondendo a 13 cursos ).
Ficaram Sem Conceito ( SC ) seis cursos ( 1,0 \% ).
A região Norte participou com 24 cursos, o correspondente a 4,0 \% do total nacional.
Desses, 12 cursos receberam o conceito 3, o que equivale a 50,0 \% do total regional.
A três cursos ( 12,5 \% ) atribuiu -se o conceito 4 ; a um curso ( 4,2 \% ), o conceito 5 ; a dois cursos ( 8,3 \% ), o conceito 1 ; e a seis ( 25,0 \% ), o conceito 2.
Nenhum curso ficou Sem Conceito ( SC ).
A região Nordeste participou com 114 cursos, o equivalente a 18,8 \% do total nacional.
Nessa região, 50 cursos ( 43,9 \% em termos regionais ) obtiveram conceito 3, o conceito modal para a região, como já comentado.
Os conceitos 4 e 5 foram atribuídos, a 12 e a 10 cursos ( 10,5 \% e 8,8 \%, respectivamente ).
Ainda nessa região, três cursos ( 2,6 \% ) foram avaliados com o conceito 1, e 38 cursos ( 33,3 \% ) receberam o conceito 2.
Apenas um curso ( 0,9 \% ) ficou Sem Conceito ( SC ).
Dos 287 cursos participantes da região Sudeste ( 47,4 \% do total nacional ), o conceito 3 foi o conceito modal, como já anteriormente assinalado, tendo sido obtido por 137 cursos ( 47,7 \% ).
O conceito 1 foi atribuído a quatro cursos ( 1,4 \% ), e o conceito 2, a 77 cursos ( 26,8 \% ).
Receberam os conceitos 4 e 5, respectivamente, 51 cursos ( 17,8 \% ) e 14 cursos ( 4,9 \% ).
Nessa região, quatro cursos ( 1,4 \% ) ficaram Sem Conceito ( SC ).
Os 119 cursos da região Sul corresponderam a 19,7 \% do total nacional.
Como já indicado, houve predominância do conceito 3, atribuído a 44,5 \% dos cursos da região ( conceito modal ), ou seja, atribuído a 53 dos 119 cursos participantes na região Sul.
O conceito 1 foi atribuído a um curso ( 0,8 \% ) ; o conceito 2, a 25 cursos ( 21,0 \% ) ; o conceito 4, a 29 cursos ( 24,4 \% ) ; e o conceito 5, a 11 cursos ( 9,2 \% ).
Nenhum curso ficou Sem Conceito.
Já dos 61 cursos participantes na região Centro-Oeste ( 10,1 \% do total nacional ), 24 ( 39,3 \% em termos regionais ) receberam conceito 2, o conceito modal.
Ao conceito 3, 20 cursos ( 32,8 \% ) foram alocados, nove cursos ( 14,8 \% ), ao conceito 4 ; quatro cursos ( 6,6 \% ), ao conceito 5, e três cursos ( 4,9 \% ), ao conceito 1.
Um curso ( 1,6 \% ) ficou Sem Conceito ( SC ).
No Gráfico 5.1b, é apresentada a distribuição cumulativa do Conceito Enade dos cursos de Bacharelado segundo Grande Região.
Em um gráfico de distribuição cumulativa, em cada valor de abscissa ( eixo x ), o valor apresentado no eixo das ordenadas ( eixo y ) é o valor acumulado da variável ( Conceito Enade ).
Observando -se a linha vertical pontilhada, que \textbf{ajuda} a visualizar o percentual de conceitos acumulados até o valor 3, ou seja, os conceitos 1, 2 e 3, pode -se notar que, por exemplo, a região Norte ( linha azul ) apresentou 83,3 \% dos cursos com conceito 3 ou abaixo desse valor.
A região Sul ( linha rosa ), por outro lado, apresentou 66,4 \% dos cursos na mesma situação : com conceito 3 ou abaixo.
Em linhas gerais, em um gráfico de distribuição cumulativa, poligonais mais à esquerda representam Grandes Regiões com pior distribuição, e poligonais mais à direita, as Grandes Regiões com os melhores conceitos.
A região Norte ( linha azul ) apresentou os cursos com os conceitos mais baixos, ao passo que a região Sul ( linha rosa ) apresentou as distribuições com os valores maiores.
Cursos em IES situadas nas regiões Sudeste ( linha vermelha ), Nordeste ( linha verde ) e Centro-Oeste ( linha \textbf{laranja} ) apresentaram uma situação intermediária. 5.2 CONCEITOS POR CATEGORIA ADMINISTRATIVA, POR MODALIDADE DE ENSINO E GRANDE REGIÃO Na Tabela 5.2a, é apresentada a distribuição dos cursos dos cursos de Licenciatura de Educação Física participantes do Enade/2021, por Categoria Administrativa e Modalidade de Ensino, de acordo com os conceitos por eles alcançados, segundo as Grandes Regiões brasileiras.
Dos 542 cursos participantes, 140 ( 25,8 \% ) eram ministrados em Instituições Públicas, e 402 ( 74,2 \% ), em Instituições Privadas.
Quanto à Modalidade de Ensino, 45 cursos eram a Distância, e 497, Presenciais.
De acordo com as informações da Tabela 5.2a, em termos nacionais, entre Instituições Públicas e Privadas, observa -se que, dos 24 cursos avaliados com conceito 5, 17 eram oferecidos em IES Públicas, e sete, em IES Privadas.
Dos 140 cursos participantes de IES Públicas, o conceito 3 foi o valor modal, atribuído a 55 cursos ( 39,3 \% ).
Dos demais cursos de Instituições Públicas participantes, 47 cursos ( 33,6 \% ) foram avaliados com conceito 4 ; 17 cursos ( 12,1 \% ), como já comentado, com conceito 5 ; um curso ( 0,7 \% ) recebeu o conceito 1 ; e 20 cursos ( 14,3 \% ) obtiveram o conceito 2.
Nessa categoria, nenhum curso ficou Sem Conceito ( SC ).
Na rede privada, o conceito modal foi 3, atribuído a 164 cursos, o correspondente a 40,8 \% dos 402 cursos da categoria.
Dos demais cursos participantes, 66 ( 16,4 \% ) receberam o conceito 4, sete ( 1,7 \% ), como já comentado, o conceito 5.
O conceito 1 foi atribuído a 12 cursos ( 3,0 \% ), e o conceito 2, a 130 cursos ( 32,3 \% ).
Nessa Categoria Administrativa, 23 cursos ( 5,7 \% ) ficaram Sem Conceito ( SC ).
Dos 497 cursos participantes da Modalidade Presencial, o conceito 3 foi o valor modal, atribuído a 201 cursos ( 40,4 \% ).
Dos demais cursos nessa Modalidade de Ensino, 110 cursos ( 22,1 \% ) receberam o conceito 4, 24 cursos ( 4,8 \% ) receberam o conceito 5, 12 cursos ( 2,4 \% ), o conceito 1, e 130 cursos ( 26,2 \% ), o conceito 2.
Nessa Modalidade de Ensino, 20 cursos ( 4,0 \% ) ficaram Sem Conceito ( SC ).
Aos 45 cursos a Distância foram atribuídos os conceitos : 2 a 20 cursos ( 44,4 \%, o conceito modal ) ; 3 a 18 cursos ( 40,0 \% ) ; 4 a três cursos ( 6,7 \% ) ; e 1 a um curso ( 2,2 \% ).
Nessa Modalidade de Ensino, três cursos ( 6,7 \% ) ficaram Sem Conceito ( SC ), e nenhum recebeu o conceito 5.
Na análise por região, observa -se que, na região Norte, as Instituições Privadas participaram com 16 cursos ( 53,3 \% do total regional ), e um curso ficou Sem Conceito ( SC ).
O conceito modal para as Instituições Privadas na região foi 3, atribuído a 7 cursos, correspondendo a 43,8 \% dessa categoria na região ; já os conceitos 4 e 5 foram atribuídos a um curso cada ( 6,3 \% em cada conceito ).
Seis cursos ( 37,5 \% ) foram avaliados com o conceito 2, e um curso ficou Sem Conceito ( SC ).
As Instituições Públicas participaram com 14 cursos na região Norte ( 46,7 \% do total regional ), aos quais foram atribuídos os conceitos : 3 a sete cursos ( 50,0 \% ), 4 a quatro cursos ( 28,6 \% ) e 2 a três cursos ( 21,4 \% ).
Nenhum curso ficou Sem Conceito ( SC ) ou recebeu os conceitos 1 ou 5, nessa combinação de categoria e região.
Na região Norte, a maioria dos cursos foi oferecida na Modalidade Presencial : 28 de 30 cursos, e 13 cursos receberam o conceito modal 3.
Os demais cursos presenciais da região receberam os conceitos 2 ( oito cursos ), 4 ( cinco cursos ) e 5 ( um curso ).
Um curso ficou Sem Conceito ( SC ), e nenhum recebeu conceito 1 na Modalidade Presencial na região Nordeste.
Dos dois cursos a Distância participantes da região Norte, um recebeu o conceito 2 ; e um, o conceito 3.
Na região Nordeste, a rede privada concentrou 55 dos 101 cursos participantes, o equivalente a 54,5 \% do total da região.
Entre os cursos oferecidos por IES Privadas no Nordeste, 25 cursos foram avaliados com conceito 3, o conceito modal.
Oito cursos obtiveram o conceito 4 ; dois receberam o conceito 5 ; dois cursos, o conceito 1 ; e 14 cursos ; o conceito 2.
Nessa combinação de categoria e região, quatro cursos ficaram Sem Conceito ( SC ).
As Instituições Públicas da região Nordeste participaram com 46 cursos ( 45,5 \% do total da região ).
Desses, 18 obtiveram conceito 4, o valor modal ; seguidos por 17 cursos, que receberam o conceito 3 ; seis que receberam o conceito 5 ; e cinco que receberam o conceito 2.
Nenhum curso ficou Sem Conceito ( SC ) ou com o conceito 1 nessa região.
Na região Nordeste, a maioria dos cursos foi oferecida na Modalidade Presencial : 96 de 101 cursos e 40 cursos presenciais receberam o conceito modal 3.
Os demais cursos Presenciais da região obtiveram os conceitos 1 ( dois cursos ), 2 ( 17 cursos ), 4 ( 25 cursos ) e 5 ( oito cursos ).
Quatro cursos ficaram Sem Conceito ( SC ) na Modalidade Presencial na região Nordeste.
Dos cinco cursos a Distância participantes da região Nordeste, dois receberam o conceito 2 ; dois, o conceito 3 ; e apenas um recebeu o conceito 4.
Na região Sudeste, a proporção de cursos da rede privada, 83,7 \%, correspondeu a 210 dos 251 cursos participantes.
Nessa categoria e região, o conceito modal foi 3, atribuído a 84 cursos, tendo 12 cursos ficado Sem Conceito ( SC ).
Os demais foram avaliados com os seguintes conceitos : 1 ( quatro cursos ) ; 2 ( 73 cursos ) ; 4 ( 33 cursos ) ; e 5 ( quatro cursos ).
Dos 41 cursos oferecidos por Instituições Públicas na região Sudeste, os conceitos 3 e 4 foram atribuídos a 14 cursos cada, uma situação bimodal.
Nessa região, oito cursos receberam o conceito 5, e cinco, o conceito 2.
Nenhum curso recebeu o conceito 1 ou ficou Sem Conceito ( SC ).
Também na região Sudeste, a quase totalidade foi oferecida na Modalidade Presencial : 234 de 251 cursos.
Para esses, a moda foi o conceito 3, com 93 cursos.
Os demais cursos presenciais da região foram alocados nos conceitos 1 ( quatro cursos ) ; 2 ( 68 cursos ) ; 4 ( 47 cursos ) ; e 5 ( 12 cursos ).
Ficaram Sem Conceito ( SC ) dez cursos.
Dos 17 cursos oferecidos na Modalidade Educação a Distância na região Sudeste, dez obtiveram o conceito 2 ( modal ), e cinco, o conceito 3.
Nenhum curso recebeu os conceitos 1, 4 ou 5.
Dois cursos ficaram Sem Conceito ( SC ).
As Instituições Privadas concentraram 87 dos 109 cursos participantes da região Sul, 79,8 \% do total regional.
Desses, 36 cursos obtiveram conceito 3, o conceito modal.
Nessa combinação de Categoria Administrativa e Grande Região, dois cursos receberam o conceito 1 ; 24 cursos, o conceito 2 ; e 21 cursos, o conceito 4.
Nenhum curso recebeu o conceito 5.
Quatro cursos ficaram Sem Conceito ( SC ).
As Instituições Públicas na região Sul participaram com 22 cursos ( 20,2 \% dos cursos da região ), aos quais foram atribuídos os conceitos 4 ( nove cursos, conceito modal ), 5 ( dois cursos ), 2 ( três cursos ) e 3 ( oito cursos ).
Nenhum curso da região Sul oferecido por IES Públicas ficou Sem Conceito ( SC ) ou recebeu o conceito 1.
Dos 109 cursos participantes da região Sul, 91 eram Presenciais.
O conceito modal para essa Modalidade de Ensino foi 3, atribuído a 36 cursos.
O conceito 1 foi atribuído a um curso ; o conceito 2, a 21 cursos ; o conceito 4, a 28 cursos ; e o conceito 5, a dois cursos.
Três cursos Presenciais da região Sul ficaram Sem Conceito ( SC ).
Já dos 18 cursos oferecidos nessa região na Modalidade a Distância, oito receberam o conceito modal 3.
Nessa Modalidade, o conceito 1 foi atribuído a um curso, o conceito 2 foi atribuído a seis cursos, e o conceito 4, a dois cursos.
Um curso ficou Sem Conceito ( SC ), e nenhum recebeu o conceito 5.
Na região Centro Oeste, 34 dos 51 cursos participantes eram de Instituições Privadas ( 66,7 \% em termos regionais ).
Desses, 13 cursos receberam o conceito 2, modal.
Dos demais cursos, 12 receberam o conceito 3, três receberam o conceito 4, e quatro, o conceito 1.
Nenhum curso recebeu o conceito 5.
Dois cursos ficaram Sem Conceito ( SC ).
Quanto aos 17 cursos oferecidos por Instituições Públicas na região Centro Oeste ( 33,3 \% do total regional ), nove cursos apresentaram conceito 3, o conceito modal, e dois receberam o conceito 4.
Os conceitos 1 e 5 foram atribuídos a um curso cada.
Nenhum curso ficou Sem Conceito ( SC ) nessa combinação de categoria e região.
Na região Centro-Oeste, participaram 48 cursos Presenciais, e o conceito modal foi 3, recebido por 19 cursos.
Quanto aos demais cursos, os conceitos 1 e 4 foram atribuídos a cinco cursos cada ; o conceito 2, a 16 cursos ; e o conceito 5, a um curso.
Nessa Modalidade, dois cursos ficaram Sem Conceito ( SC ).
Dos três cursos a Distância da região Centro-Oeste, dois receberam o conceito 3, e um recebeu o conceito 2.
No Gráfico 5.2a, apresenta -se a distribuição do Conceito Enade dos cursos de Licenciatura, segundo a Categoria Administrativa da IES.
Os cursos em IES Públicas ( linha azul ) apresentam uma poligonal mais à direita do que os em IES Privadas ( linha verde ) e, consequentemente, uma distribuição de Conceito Enade com valores maiores.
Na Tabela 5.2b, é apresentada a distribuição dos cursos dos cursos de Bacharelado de Educação Física participantes do Enade/2021, por Categoria Administrativa e Modalidade de Ensino, de acordo com os conceitos por eles alcançados, segundo as Grandes Regiões brasileiras.
Dos 605 cursos participantes, 78 ( 12,9 \% ) eram ministrados em Instituições Públicas, e 527 ( 87,1 \% ), em Instituições Privadas.
Quanto à Modalidade de Ensino, 28 cursos eram a Distância, e 577, Presenciais.
De acordo com as informações da Tabela 5.2b, em termos nacionais, entre Instituições Públicas e Privadas, observa -se que, dos 40 cursos avaliados com o conceito 5, 24 eram oferecidos em IES Públicas, e 16, em IES Privadas.
Dos 78 cursos participantes de IES Públicas, o conceito 5 foi o valor modal, atribuído a 24 cursos ( 30,8 \% ).
Dos demais cursos de Instituições Públicas participantes, um curso ( 1,3 \% ) foi avaliado com o conceito 1 ; dez cursos ( 12,8 \% ) foram avaliados com o conceito 2 ; 20 cursos ( 25,6 \% ), com o conceito 3 ; e 23 cursos ( 29,5 \% ), com o conceito 4.
Nenhum curso ficou Sem Conceito ( SC ).
Na rede privada, o conceito modal foi 3, atribuído a 252 cursos, o correspondente a 47,8 \% dos 527 cursos da categoria.
Dos demais cursos participantes, 12 ( 2,3 \% ) receberam o conceito 1 ; 160 ( 30,4 \% ), o conceito 2 ; 81 ( 15,4 \% ), o conceito 4 ; e 16 ( 3,0 \% ), o conceito 5.
Nessa Categoria Administrativa, seis cursos ( 1,1 \% ) ficaram Sem Conceito ( SC ).
Dos 577 cursos participantes na Modalidade Presencial, o conceito 3 foi o valor modal, atribuído a 264 cursos ( 45,8 \% ).
Dos demais cursos nessa Modalidade de Ensino, 11 cursos ( 1,9 \% ) receberam o conceito 1 ; 158 cursos ( 27,4 \% ), o conceito 2 ; 100 cursos ( 17,3 \% ), o conceito 4 ; e 39 cursos ( 6,8 \% ), o conceito 5.
Nessa Modalidade de Ensino, cinco cursos ( 0,9 \% ) ficaram Sem Conceito ( SC ).
Aos 28 cursos a Distância informados, atribuíram -se os seguintes conceitos : 1 a dois cursos ( 7,1 \% ) ; 2 a 12 cursos ( 42,9 \% ), o conceito modal ; 3 a oito cursos ( 28,6 \% ) ; 4 a quatro cursos ( 14,3 \% ) ; e 5 a apenas um curso ( 3,6 \% ).
Nessa Modalidade de Ensino, apenas um curso ( 3,6 \% ) ficou Sem Conceito ( SC ).
Na análise por região, observa -se que, na região Norte, todos os 24 cursos oferecidos na região Norte foram Presenciais.
As Instituições Privadas participaram com 20 cursos ( 83,3 \% do total regional ).
O conceito modal para as Instituições Privadas na região foi 3, atribuído a nove cursos, correspondendo a 45,0 \% dessa categoria na região.
Um curso foi avaliado com o conceito 1 ; seis cursos foram avaliados com o conceito 2 ; três cursos, com o conceito 4 ; e um recebeu o conceito 5.
As Instituições Públicas participaram com quatro cursos na região Norte ( 16,7 \% do total regional ).
Três cursos receberam o conceito 3, e um recebeu o conceito 1.
Nenhum curso na região Norte ficou Sem Conceito ( SC ).
Na região Nordeste, a rede privada concentrou 95 dos 114 cursos participantes, o equivalente a 83,3 \% do total da região.
Entre os cursos oferecidos por IES Privadas no Nordeste, 46 cursos foram avaliados com conceito 3, o conceito modal.
Três cursos obtiveram o conceito 1 ; 33 cursos, o conceito 2 ; oito cursos, o conceito 4 ; e quatro cursos, o conceito 5.
Nessa combinação de categoria e região, um curso ficou Sem Conceito ( SC ).
As Instituições Públicas da região Nordeste participaram com 19 cursos ( 16,7 \% do total da região ).
Desses, seis cursos obtiveram conceito 5, o valor modal ; já os conceitos 3 e 4 foram atribuídos a quatro cursos cada.
Cinco cursos receberam o conceito 2.
Nenhum curso ficou Sem Conceito ( SC ) ou recebeu o conceito 1.
Na região Nordeste, a maioria dos cursos foi oferecida na Modalidade Presencial : 111 de 114 cursos, sendo que 49 cursos receberam o conceito modal 3.
Os demais cursos Presenciais da região receberam os seguintes conceitos : 1 ( três cursos ), 2 ( 36 cursos ), 4 ( 12 cursos ) e 5 ( dez cursos ).
Um curso ficou Sem Conceito ( SC ) na Modalidade Presencial na região Nordeste.
Dos três cursos a Distância participantes da região Nordeste, dois receberam o conceito 2, e um recebeu o conceito 3.
Na região Sudeste, a proporção de cursos da rede privada, 88,5 \%, correspondeu a 254 dos 287 cursos participantes.
Nessa categoria e região, o conceito modal foi 3, atribuído a 128 cursos, e quatro cursos ficaram Sem Conceito ( SC ).
Os demais foram avaliados com os seguintes conceitos : 1 ( quatro cursos ) ; 2 ( 74 cursos ) ; 4 ( 41 cursos ) ; e 5 ( três cursos ).
Dos 33 cursos oferecidos por Instituições Públicas na região Sudeste, o conceito modal foi 5, atribuído a 11 cursos.
Nessa região, três cursos de IES Públicas receberam o conceito 2 ; nove cursos, o conceito 3 ; e dez, o conceito 4.
Nenhum curso ficou Sem Conceito ( SC ) ou foi avaliado com o conceito 1.
Também na região Sudeste, a maioria dos cursos foi oferecida na Modalidade Presencial : 273 cursos.
Para esses, a moda foi o conceito 3, com 134 cursos.
Os demais cursos Presenciais da região receberam os conceitos : 1 ( dois cursos ), 2 ( 72 cursos ), 4 ( 48 cursos ) e 5 ( 14 cursos ).
Ficaram Sem Conceito ( SC ) três cursos.
Na região Sudeste, os 14 cursos oferecidos na Modalidade Educação a Distância obtiveram os conceitos : 2 ( cinco cursos, modal ), 3 e 4 ( três cursos cada ) e 1 ( dois cursos ).
Nenhum curso obteve o conceito 5, e um curso ficou Sem Conceito ( SC ).
As Instituições Privadas concentraram 106 dos 119 cursos participantes da região Sul, 89,1 \% do total regional.
Desses, 52 cursos obtiveram conceito 3, o conceito modal.
Nessa combinação de Categoria Administrativa e Grande Região, um curso recebeu conceito 1 ; 25 cursos receberam o conceito 2 ; 22 cursos, o conceito 4 ; e seis cursos, o conceito 5.
Nenhum curso ficou Sem Conceito ( SC ).
As Instituições Públicas na região Sul participaram com 13 cursos ( 10,9 \% dos cursos da região ) : sete cursos receberam o conceito 4 ( modal ), cinco cursos receberam o conceito 5, e um recebeu o conceito 3.
Nenhum curso ficou Sem Conceito ( SC ) ou recebeu os demais conceitos.
Dos 119 cursos participantes da região Sul, 111 eram Presenciais.
O conceito modal para essa Modalidade de Ensino foi o 3, atribuído a 50 cursos.
O conceito 1 foi atribuído a um curso ; o conceito 2, a 22 cursos ; o conceito 4, a 28 cursos ; e o conceito 5, a dez cursos.
Nenhum curso Presencial da região Sul ficou Sem Conceito ( SC ).
Já dos oito cursos oferecidos nessa região na Modalidade a Distância, três receberam o conceito 2, e três, o conceito 3 ( bimodal ) ; já os conceitos 4 e 5 foram atribuídos a um curso cada.
Nenhum curso ficou Sem Conceito ( SC ) ou recebeu o conceito 1.
Na região Centro-Oeste, 52 dos 61 cursos participantes eram de Instituições Privadas ( 85,2 \% em termos regionais ).
Desses, 22 receberam conceito 2, o conceito modal.
Dos demais cursos, três receberam o conceito 1 ; 17, o conceito 3 ; sete, o conceito 4 ; e dois, o conceito 5.
Um curso ficou Sem Conceito ( SC ).
Quanto aos nove cursos oferecidos por Instituições Públicas na região Centro-Oeste ( 14,8 \% do total regional ), três cursos apresentaram conceito 3, o conceito modal.
Os conceitos 2, 4 e 5 foram atribuídos a dois cursos cada.
Nenhum curso ficou Sem Conceito ( SC ) ou recebeu o conceito 1 nessa combinação de categoria e região.
Na região Centro-Oeste, participaram 58 cursos Presenciais, e o conceito modal foi o 2, recebido por 22 cursos.
Quanto aos demais cursos, o conceito 1 foi atribuído a três cursos ; o conceito 3, a 19 cursos ; o conceito 4, a nove cursos ; e o conceito 5, a quatro cursos.
Nessa Modalidade, um curso ficou Sem Conceito ( SC ).
Dos três cursos a Distância da região Centro-Oeste, dois receberam o conceito 2, e um recebeu o conceito 3.
No Gráfico 5.2b, apresenta -se a distribuição do Conceito Enade dos cursos de Bacharelado, segundo a Categoria Administrativa da IES.
Os cursos em IES Públicas ( linha azul ) apresentam uma poligonal mais à direita do que os em IES Privadas ( linha verde ) e, consequentemente, uma distribuição de Conceitos Enade com valores maiores. 5.3 CONCEITOS POR ORGANIZAÇÃO ACADÊMICA E GRANDE REGIÃO Na Tabela 5.3a, encontra -se a distribuição dos conceitos atribuídos aos cursos de Licenciatura participantes do Enade/2021, na Área de Educação Física, por Organização Acadêmica, segundo as Grandes Regiões brasileiras.
Dos 542 cursos de Educação Física participantes, 255 eram oferecidos em Universidades, 148 em Centros Universitários, 128 em Faculdades e 11 em CEFET/IF.
Esta distribuição corresponde a, respectivamente, 47,0 \%, 27,3 \%, 23,6 \% e 2,0 \% dos cursos.
De acordo com os dados apresentados, de todos os 24 cursos avaliados com o conceito 5, 19 eram vinculados a Universidades.
Os cursos vinculados a esse tipo de Organização Acadêmica tiveram o conceito 3 como conceito modal, atribuído a 107 cursos ( 42,0 \% ).
Os demais cursos vinculados a Universidades receberam os conceitos 1 ( três cursos ), 2 ( 57 cursos ), 4 ( 67 cursos ) e 5 ( 19 cursos, como já mencionado ).
Dois cursos ficaram Sem Conceito ( SC ).
Dos cursos em Centros Universitários, o conceito modal também foi o 3, atribuído a 66 cursos ( 44,6 \% ).
Dez cursos vinculado a esse tipo de Organização Acadêmica ficaram Sem Conceito ( SC ), e os outros cursos receberam os conceitos 1 ( quatro cursos ), 2 ( 39 cursos ), 4 ( 26 cursos ) e 5 ( três cursos ).
Dos 128 cursos mantidos por Faculdades, 53 ( 41,4 \% ) receberam conceito 2, o conceito modal.
O segundo conceito mais frequentemente atribuído aos cursos oferecidos por essa categoria de IES foi o conceito 3, recebido por 41 cursos.
Na sequência, 15 cursos obtiveram o conceito 4 ; seis cursos, o conceito 1 ; e dois cursos, o conceito 5.
Dos cursos oferecidos em Faculdades, onze ficaram Sem Conceito ( SC ).
Na Organização Acadêmica CEFET/IF, dos 11 cursos que participaram, os conceitos 3 e 4 foram atribuídos a cinco cursos cada, e um curso recebeu o conceito 2.
Nenhum curso ficou Sem Conceito ( SC ) ou recebeu os demais conceitos.
Considerando -se, separadamente, as regiões brasileiras, verifica -se que, na região Norte, as Universidades participaram com 13 dos 30 cursos avaliados, aos quais foi atribuído o conceito 3 para seis cursos, conceito modal, o 2 para quatro cursos, e o 4 para três cursos.
Nenhum curso ficou Sem Conceito ( SC ) ou recebeu os demais conceitos nessa combinação de categoria e região.
Os Centros Universitários da região Norte foram representados por oito cursos, dos quais quatro obtiveram o conceito 2, modal.
Os demais cursos foram avaliados com os conceitos 3 ( três cursos ) e 5 ( um curso ).
Nenhum curso ficou Sem Conceito ( SC ) ou recebeu os demais conceitos nessa combinação de categoria e região.
As Faculdades participaram com sete cursos na região Norte, sendo que o conceito 3 ( modal ) foi atribuído a quatro cursos, e os conceitos 2 e 4 foram atribuídos a um curso cada.
Um curso ficou Sem Conceito ( SC ), e nenhum curso oferecido em Faculdades da região obteve os conceitos 1 ou 5.
Dois cursos foram oferecidos em CEFET/IF aos quais foram atribuídos os conceitos 3 e 4.
Na região Nordeste, as Universidades participaram com 44 dos 101 cursos da Área de Educação Física ( Licenciatura ) oferecidos.
O conceito modal 4 foi atribuído a 19 cursos.
Cinco cursos obtiveram o conceito 2 ; 14 cursos, o conceito 3 ; e seis, o conceito 5.
Nenhum curso oferecido por Universidades no Nordeste recebeu o conceito 1 ou ficou Sem Conceito ( SC ).
Os Centros Universitários contaram com 23 cursos participantes na região Nordeste, 15 dos quais receberam o conceito modal 3.
Para os demais cursos, foram atribuídos os conceitos 2 e 4 ( três cursos cada ) e 5 ( um curso ).
Um curso oferecido por Centros Universitários no Nordeste ficou Sem Conceito ( SC ).
As Faculdades foram representadas por 30 cursos na região Nordeste, com os conceitos 2 e 3 sendo atribuídos a 11 cursos cada ( bimodal ).
Cada um dos conceitos 1 e 4 foi atribuído a dois cursos, e o conceito 5 foi atribuído a um curso.
Na região Nordeste, três cursos, mantidos por Faculdades, ficaram Sem Conceito ( SC ).
Quatro cursos foram oferecidos em CEFET/IF na região Nordeste, aos quais foram atribuídos os conceitos 3 e 4 ( dois cursos em cada conceito ).
Na região Sudeste, as Universidades concentraram 110 dos 251 cursos de Educação Física ( Licenciatura ) da região.
Dos cursos oferecidos em Universidades, nessa região, o conceito modal foi 3, atribuído a 45 cursos.
Os demais cursos receberam os conceitos 1 ( um curso ), 2 ( 30 cursos ), 4 ( 23 cursos ) e 5 ( dez cursos ).
Um curso ficou Sem Conceito ( SC ).
Os Centros Universitários participaram com 71 cursos na região Sudeste, dos quais 28 obtiveram o conceito modal 3.
O conceito 1 foi atribuído a dois cursos ; o conceito 2, a 22 cursos ; o conceito 4, a 14 curso ; e o conceito 5, a um curso.
Quatro cursos ficaram Sem Conceito ( SC ) nessa combinação de organização acadêmica e região.
Já as Faculdades foram representadas por 66 cursos na região Sudeste, dos quais sete ficaram Sem Conceito ( SC ).
O conceito modal 2 foi atribuído a 26 cursos.
Aos demais cursos foram atribuídos os conceitos 1 e 5 ( um curso cada ), 3 ( 23 cursos ) e 4 ( oito cursos ).
Quatro cursos foram oferecidos em CEFET/IF na região Sudeste, aos quais foram atribuídos os conceitos 3 e 4 ( dois cursos em cada conceito ).
Dos 109 cursos da região Sul, 65 eram mantidos por Universidades, para os quais o conceito modal foi o 3, atribuído a 31 cursos.
Os demais cursos receberam os conceitos 2 ( onze cursos ), 4 ( 20 cursos ) e 5 ( dois cursos ).
Nesse tipo de organização, um curso da região Sul ficou Sem Conceito ( SC ), e nenhum recebeu o conceito 1.
Dos 30 cursos participantes de Centros Universitários da região Sul, três ficaram Sem Conceito ( SC ), e nenhum recebeu o conceito 5.
Os conceitos restantes foram atribuídos a um curso ( conceito 1 ), nove cursos ( conceito 2 ), 11 cursos ( conceito 3, modal ) e seis cursos ( conceito 4 ).
Quanto aos 13 cursos vinculados a Faculdades na região Sul, seis receberam o conceito modal 2, um recebeu o conceito 1, dois receberam o conceito 3, e quatro, o conceito 4.
Nenhum curso ficou Sem Conceito ( SC ) ou recebeu o conceito 5.
Na região Sul, participou um curso vinculado a CEFET/IF e a ele foi atribuído o conceito 2.
Na região Centro-Oeste, 23 dos 51 cursos oferecidos eram mantidos por Universidades.
Para esse tipo de Organização Acadêmica, o conceito modal 3 foi atribuído a 11 cursos.
Os conceitos 1 e 4 foram atribuídos a dois cursos cada ; o conceito 2, a sete cursos ; e o conceito 5 foi atribuído a um curso.
Nesse tipo de organização, nenhum curso ficou Sem Conceito ( SC ).
Os Centros Universitários da região Centro-Oeste contaram com 16 cursos, dos quais nove foram avaliados com o conceito 3 ( modal ) ; três, com o conceito 4 ; um curso foi avaliado com o conceito 1 ; e um curso, com o conceito 2.
Dois cursos ficaram Sem Conceito ( SC ), e nenhum recebeu o conceito 5.
Dos doze cursos oferecidos em Faculdades na região Centro-Oeste, nove receberam o conceito modal 2, um recebeu o conceito 3, e dois receberam o conceito 1.
Nessa região, nenhum curso oferecido em Faculdades ficou Sem Conceito ( SC ) ou obteve os demais conceitos.
Não houve, na região Centro-Oeste, cursos mantidos por CEFET/IF.
No Gráfico 5.3a, apresenta -se a distribuição do Conceito Enade dos cursos de Licenciatura, segundo a Organização Acadêmica da IES.
Os cursos em Universidades ( linha azul ) apresentaram uma poligonal mais à direita do que as demais, denotando uma melhor distribuição de Conceitos Enade.
No outro extremo, encontraram -se os cursos mantidos por Faculdades ( linha vermelha ), apresentando a respectiva poligonal mais à esquerda e, consequentemente, uma distribuição de Conceitos Enade com menores valores.
Os cursos mantidos por Centros Universitários ( linha verde ) apresentaram poligonal mais central e, consequentemente, uma distribuição de Conceitos Enade com valores intermediários.
Na Tabela 5.3b, encontra -se a distribuição dos conceitos atribuídos aos cursos de Bacharelado participantes do Enade/2021, na Área de Educação Física, por Organização Acadêmica, segundo as Grandes Regiões brasileiras.
Dos 605 cursos de Educação Física participantes, 222 foram oferecidos em Universidades, 212 em Centros Universitários, 170 em Faculdades, e apenas um foi oferecido em CEFET/IF.
Essa distribuição corresponde a, respectivamente, 36,7 \%, 35,0 \%, 28,1 \% e 0,2 \% dos cursos.
De acordo com os dados apresentados, de todos os 40 cursos avaliados com o conceito 5, 29 eram vinculados a Universidades.
Os cursos vinculados a esse tipo de Organização Acadêmica tiveram o conceito 3 como conceito modal, atribuído a 95 cursos ( 42,8 \% ).
Os demais cursos vinculados a Universidades receberam os conceitos : 1 ( três cursos ), 2 ( 42 cursos ), 4 ( 51 cursos ) e 5 ( 29 cursos, como já mencionado ).
Ficaram Sem Conceito ( SC ) dois cursos.
Dos cursos em Centros Universitários, o conceito modal também foi o 3, atribuído a 105 dos 212 cursos ( 49,5 \% ).
Dois cursos vinculados a esse tipo de Organização Acadêmica ficaram Sem Conceito ( SC ) ; já os outros cursos receberam os conceitos : 1 ( três cursos ), 2 ( 60 cursos ), 4 ( 34 cursos ) e 5 ( oito cursos ).
Dos 170 cursos mantidos por Faculdades, 72 ( 42,4 \% ) receberam conceito 3, o conceito modal.
O segundo conceito mais frequentemente atribuído aos cursos oferecidos por essa categoria de IES foi o conceito 2, recebido por 68 cursos.
O conceito 1 foi atribuído a sete cursos ; o conceito 4, a 18 cursos ; e o conceito 5, a três cursos.
Dois cursos ficaram Sem Conceito ( SC ).
O único curso mantido por CEFET/IF recebeu o conceito 4.
Considerando -se, separadamente, as regiões brasileiras, verifica -se que, na região Norte, as Universidades participaram com cinco dos 24 cursos avaliados.
O conceito modal foi o 3, atribuído a quatro cursos.
O conceito 1 foi atribuído a um curso.
Nenhum curso ficou Sem Conceito ( SC ) ou recebeu os demais conceitos.
Os Centros Universitários da região Norte foram representados por 12 cursos, os quais obtiveram os conceitos 2 e 3 ( quatro cursos cada, distribuição bimodal ), 1 e 5 ( um curso cada ) e 4 ( dois cursos ).
Nenhum curso ficou Sem Conceito ( SC ).
As Faculdades participaram com sete cursos na região Norte, os quais receberam os conceitos 3 ( quatro cursos, conceito modal ), 4 ( um curso ) e 2 ( dois cursos ).
Nenhum curso recebeu os demais conceitos ou ficou Sem Conceito ( SC ).
Já os CEFET/IF não tiveram cursos avaliados na região Norte.
Na região Nordeste, as Universidades participaram com 25 dos 114 cursos da Área de Educação Física ( Bacharelado ) oferecidos.
O conceito modal foi o 3, com nove cursos ; seguido do conceito 5, com seis cursos ; e do conceito 4, com cinco cursos.
O conceito 1 foi atribuído a um curso, e o conceito 2, a três.
Um curso ficou Sem Conceito ( SC ).
Os Centros Universitários contaram com 49 cursos participantes na região Nordeste, 25 dos quais receberam o conceito modal 3.
Para os demais cursos, foram atribuídos os conceitos 2 ( 17 cursos ), 4 ( cinco cursos ) e 5 ( dois cursos ).
Nenhum curso oferecido por Centros Universitários no Nordeste ficou Sem Conceito ( SC ) ou recebeu o conceito 1.
As Faculdades foram representadas por 40 cursos na região Nordeste, 18 dos quais receberam conceito 2, o conceito modal.
Cada um dos conceitos 1, 4 e 5 foi atribuído a dois cursos.
O conceito 4 foi recebido por 16 cursos.
Nenhum curso ficou Sem Conceito ( SC ).
Os CEFET/IF não tiveram curso participante na região Nordeste.
Na região Sudeste, as Universidades concentraram 115 dos 287 cursos de Educação Física ( Bacharelado ) da região.
Dos cursos oferecidos em Universidades, nessa região, o conceito modal foi o 3, atribuído a 51 cursos.
Os demais cursos receberam os conceitos 1 ( um curso ), 2 ( 25 cursos ), 4 ( 23 cursos ) e 5 ( 14 cursos ).
Um curso ficou Sem Conceito ( SC ).
Os Centros Universitários participaram com 97 cursos na região Sudeste, dos quais 49 obtiveram o conceito modal 3.
O conceito 1 foi atribuído a dois cursos ; o conceito 2, a 26 cursos ; e o conceito 4, a 18 cursos.
Nenhum curso recebeu o conceito 5.
Dois cursos ficaram Sem Conceito ( SC ) nessa combinação de organização acadêmica e região.
Já as Faculdades foram representadas por 74 cursos na região Sudeste, dos quais um ficou Sem Conceito ( SC ).
O conceito modal 3 foi atribuído a 37 cursos.
Aos demais cursos foram atribuídos os conceitos 1 ( um curso ), 2 ( 26 cursos ) e 4 ( nove cursos ).
Nenhum curso recebeu conceito 5.
Os CEFET/IF participaram com um curso, que recebeu o conceito 4.
Dos 119 cursos da região Sul, 63 eram mantidos por Universidades, para os quais o conceito modal foi 3, atribuído a 28 cursos.
O segundo conceito mais frequente foi o 4, com 19 cursos.
O conceito 2 foi atribuído a nove cursos, e o conceito 5, a sete cursos.
Nenhum curso recebeu o conceito 1 ou ficou Sem Conceito ( SC ).
Dos 32 cursos participantes de Centros Universitários da região Sul, 17 cursos foram avaliados com o conceito 3 ( conceito modal ) ; sete cursos foram avaliados com o conceito 2 ; cinco cursos, com o conceito 4 ; e três cursos, com o conceito 5.
Nenhum curso recebeu o conceito 1 ou ficou Sem Conceito ( SC ).
Quanto aos 24 cursos vinculados a Faculdades na região Sul, nove receberam o conceito modal 2.
O conceito 1 foi atribuído a um curso ; o 3, a oito ; o 4, a cinco ; e o 5, a um curso.
Nenhum curso ficou Sem Conceito ( SC ).
Já os CEFET/IF não participaram na avaliação de cursos na região Sul.
Na região Centro-Oeste, 14 dos 61 cursos oferecidos eram mantidos por Universidades.
Para esse tipo de Organização Acadêmica, nenhum curso ficou Sem Conceito ( SC ) ou recebeu o conceito 1.
O conceito modal foi o 2 ( cinco cursos ), seguido dos conceitos 4 ( quatro cursos ), 3 ( três cursos ) e 5 ( dois cursos ).
Os Centros Universitários da região Centro-Oeste contaram com 22 cursos, que foram avaliados com os conceitos : 3 ( dez cursos, modal ), 2 ( seis cursos ), 4 ( quatro cursos ) e 5 ( dois cursos ).
Nenhum curso recebeu o conceito 1 ou ficou Sem Conceito ( SC ).
As Faculdades participaram com 25 cursos, dos quais um ficou sem conceito.
Dos demais, 13 receberam o conceito modal 2 ; já os demais conceitos foram : 1 ( três cursos ), 3 ( sete cursos ) e 4 ( um curso ).
Não houve curso avaliado como conceito 5 ; já os CEFET/IF não tiveram cursos participantes na região Centro-Oeste.
No Gráfico 5.3b, apresenta -se a distribuição do Conceito Enade dos cursos de Bacharelado, segundo a Organização Acadêmica da IES.
Nos CEFET/IF ( em \textbf{laranja} ), o único curso dessa Organização Acadêmica recebeu o conceito 4, então a poligonal cumulativa alcançou 100 \% para o conceito 4.
Os cursos em Universidades ( linha azul ) apresentaram uma poligonal mais à direita do que as demais, denotando uma melhor distribuição de Conceitos Enade.
No outro extremo, com uma distribuição muito semelhante entre si, encontram -se os cursos mantidos por Centros Universitários ( linha verde ) e por Faculdades ( linha vermelha ), apresentando as respectivas poligonais mais à esquerda e, consequentemente, uma distribuição de Conceitos Enade com menores valores.

% matched lemmas: ajudar, laranja
\end{textsample}
